\chapter{28}



Das Gebiet präsentierte sich wie eine neu errichtete Theaterkulisse, 
eben bemalt, lackiert und zum Trocknen gestellt. Und er als einziger
mittendrin. Einige Meter abseits des Wanderpfades entdeckte er eine
kleine, nahe beisammenstehende Gruppe Tannen. Er legte sich darunter und döste
augenblicklich weg.

Es war ein Sonnenstrahl, der seine rechte Gesichtshälfte erhitzte und
ihn aufwachen liess. Noch mit geschlossenen Augen wischte er mit der
Handfläche über seine Wange und blieb vorerst liegen. Sein Körper
fühlte sich schwer an.
Als er die Augen schmal öffnete, versank sein Blick in einer Ansammlung
von Grüntönen, einem Blättermeer, das sich dicht über ihm zu einem
Dach formierte. Das Bild glich einem Gemälde, lose aufgehängt unter dem 
Himmel. Der späte Nachmittag zeigte sich überaus versöhnlich. Louis setzte sich auf. Im
Unterschied zum Tal liess sich hier oben auffallend leicht atmen. Keine
Abgase. Autofreie Zone. Beeindruckend. Wo gab es so was noch.
Beim Blick auf sein Handy wunderte er sich. Bereits so spät? Er hatte
lange geschlafen. Hungrig öffnete er seinen Rucksack, holte Brot und
Oliven heraus und schluckte die ersten Bissen hastig hinunter.
Erstaunlich, wie gut ihm diese einfachen Nahrungsmittel schmeckten. 

Mit der Aussicht auf die ersehnte Intimität mit seiner \emph{Gibson}
schlug sein Herz höher. Gleich nur er und sie. Er konnte es kaum glauben
und öffnete aufgeregt den Gitarrenkoffer.
Das Instrument war verstimmt. Sorgfältig begann er an den Stimmwirbeln
zu drehen und die entsprechenden Saiten anzuschlagen. 
Wie sehr sich Louis vor ein paar Tagen noch danach gesehnt hatte, seine
neue Gitarre mit «L’amore è» einzuweihen. In Erinnerung an
Giorgia und ihn und an glückliche Stunden. Doch es fühlte sich falsch
an. Passte nicht. Weder hier hin, noch zu Ciro.
Louis rückte sich zurecht. Als er das Instrument startbereit in den Armen hielt,
überkam ihn ein tiefes Gefühl von Dankbarkeit. Er blinzelte einzelne
Tränen weg und dachte dabei an Meo und Vanessa. Sie hatten ihm Glück
gebracht.
Die ersten Akkorde liess er andächtig ausklingen. Er kam nach und nach in
Fahrt und ohne nachzudenken, fand er sich mittendrin. Klänge und
Rhythmen lösten sich ab. Neue Melodien entstanden. Zuweilen summte er
mit. Louis wähnte sich in Zeitlosigkeit, angenehm aufgehoben, und spielte lange.
Bis ihn die Veränderung des Lichts an die Zeit erinnerte. Bald siebzehn
Uhr. Louis trank die Flasche leer.
Laut Wegbeschreibung musste der Hof ganz in der Nähe liegen. Er schätzte
die Entfernung auf eine Viertelstunde Gehzeit ein.

Er verstaute die Einkaufstüte im Rucksack, schwang diesen auf den Rücken und nahm den Gitarrenkoffer hoch. 
Mit gelegentlichem Blick auf sein Handy folgte er dem Weg. Es ging
steiler hoch als bis anhin. Wie leicht er vorankam. Die Pause unter den
Bäumen hatte seine Batterien neu aufgeladen.
